\documentclass[11pt]{article}
\usepackage[margin=1in]{geometry}
\usepackage{booktabs}
\usepackage{url}
\title{The estimator's ego bound is lost during avoidance, and a correction}
\date{Prepared October 4, 2026}
\begin{document}
\maketitle

\section{Question}
The user asked for the controller's observer margins to use the estimator's
design-guaranteed error bounds at each future frame, instead of an open-loop
tube. Before implementing this, the published bounds were inspected in a noisy
15-m/s headOn run. At 1.16~s the estimator had already declared its ego bound
invalid: \texttt{egoValid} false, yaw bound infinite, every published ego and
target enclosure unavailable. The reason was
\texttt{ego-measurement-inconsistent-with-declared-bounds}.

\section{Replay}
\path{scripts/replayEstimatorValidity.m} replays the NRMM adapter on the
recorded truth of each noisy campaign run, with the same configuration and
seed. The adapter is deterministic, so this reproduces the published estimate
without calling the controller. Results for the current source (aec20bc,
\path{~/.cache/collisionAvoidance/clf-time-constant-20261004/noisy}):

\begin{center}
\begin{tabular}{rlrrr}
\toprule
speed & scenario & ego bound invalid from & $|r|$ there & peak $|r|$\\
\midrule
8 & headOn & 0.35 s & 0.758 & 0.939\\
8 & acceleratingHeadOn & 0.35 s & 0.759 & 0.936\\
8 & brakingLead & 1.50 s & 0.742 & 0.742\\
8 & crossing & 0.50 s & 0.727 & 0.986\\
8 & turningCrossing & 0.40 s & 0.778 & 0.792\\
8 & curvedHeadOn & 0.35 s & 0.765 & 0.944\\
8 & curvedCrossing & no frames & & \\
15 & headOn & 0.85 s & 0.753 & 0.990\\
15 & acceleratingHeadOn & valid (15 frames) & & 0.690\\
15 & brakingLead & valid (233 frames) & & 0.694\\
15 & crossing & 1.70 s & 0.778 & 0.890\\
15 & turningCrossing & valid (17 frames) & & 0.661\\
15 & curvedHeadOn & 1.30 s & 0.788 & 0.811\\
15 & curvedCrossing & valid (20 frames) & & 0.677\\
\bottomrule
\end{tabular}
\end{center}
In 9 of the 13 runs with frames, the avoidance maneuver drives the yaw rate
above the estimator's declared domain,
\texttt{observer.ego.domain.yawRateMaximum} $=0.75$~rad/s (plus gyro noise).
The bound is then invalidated and stays invalid for the rest of the run. From
that point the controller receives no uncertainty information and solves the
certainty-equivalent problem. In that problem the prefix collision rows are
soft, at the user's earlier choice. The yaw rates there (0.73--0.79~rad/s
falling, up to 0.99~rad/s peak) are within the plant's capability; the declared
estimator domain, not the vehicle, is exceeded. In the 8-m/s crossing run the
speed also fell to 4.79~m/s, below the declared ego speed floor of 5~m/s,
after the bound was already invalid.

\section{Correction of an earlier explanation}
\path{LQR_REMOVAL_20261004.tex} attributed the noisy 15-m/s brakingLead
collision to an open-loop ego tube combined with soft observer tightening.
That is not supported. A replay of that run (8ae93c3,
\path{no-lqr-20261004b/noisy}) shows the ego bound invalid from 1.00~s
($|r|=0.779$~rad/s). The collision was at 9.21~s, and the slack cap stayed at
the certainty-equivalent tolerance $10^{-4}$ until 9.10~s. Positive slack
appeared only at 9.15 and 9.20~s, on the soft nominal prefix rows. During
the last 8~s no observer tightening or ego tube was in the problem.

The collision followed from a wrong target estimate (acceleration $+0.1$
against a true $-0.5$~m/s$^2$; about 0.47~m of position error) with only the
0.2-m nominal margin and soft prefix collision rows. The two counterfactuals of
that report differ in their first second, while the bound was still valid.
They therefore changed the later trajectory, but they do not identify the tube
as the mechanism. The closing remark of \path{CLF_TIME_CONSTANT_20261004.tex},
which repeats that mechanism, is corrected by this report as well.

\section{Consequence for the requested change}
Future bounds from the estimator's comparison system are well defined only
while the estimator's premises hold. Three channels have no such bound in the
present configuration:
\begin{itemize}
\item \emph{Ego velocity.} The ego acceleration bound is declared infinite
  (\texttt{accelerationNormMaximum}), so the velocity comparison is not used
  between samples. Only the measured residual of each sample bounds velocity
  error, and nothing bounds it a priori.
\item \emph{Yaw.} The yaw bound has no guaranteed contraction; it grows with
  gyro noise unless the course correspondence is informative.
\item \emph{After a domain violation.} Once the ego domain is violated, no
  bound exists at all.
\end{itemize}
The future-bound tube was therefore not implemented. Two coupled designs
would restore valid bounds: the controller keeps the vehicle inside the
estimator's declared domain (a yaw-rate and speed constraint in the plan),
or the estimator's domain and acceleration bound are enlarged to the
vehicle's physical limits and its gains re-synthesized.
\end{document}
